%% chapter.tex — Chapter template for documentation-pro projects
%% Copy to chapters/chNN-title.tex and fill in content

\chapter{Chapter Title}
\label{ch:chapter-label}

%% SECTION 1: OVERVIEW (150–300 words)
\section{Overview}
\label{sec:chapter-label-overview}

\chapterlead{%
  [One to three motivating sentences. What does this chapter cover?
  Why does it matter? What will the reader be able to do?]
}

By the end of this chapter, you will be able to:
\begin{itemize}
  \item [Skill 1]
  \item [Skill 2]
  \item [Skill 3]
\end{itemize}

%% SECTION 2: MOTIVATION (200–400 words)
\section{Motivation}
\label{sec:chapter-label-motivation}

[Present the problem with a concrete scenario. Do not start with the solution.]

%% SECTION 4: CORE CONCEPTS (500–1500 words)
\section{Core concepts}
\label{sec:chapter-label-concepts}

[Define every major concept. Follow each abstraction with a concrete example.]

\subsection{First concept}
\label{subsec:chapter-label-concept1}

[Definition, motivation, concrete example, diagram if applicable.]

\begin{figure}[H]
  \centering
  \input{figures/chNN/concept-diagram}
  \caption{[What is shown and why it matters.]}
  \label{fig:chapter-label-concept1}
\end{figure}

%% SECTION 5: ARCHITECTURE
\section{Architecture}
\label{sec:chapter-label-architecture}

[Structural overview with at least one architecture diagram.]

%% SECTION 6: IMPLEMENTATION (500–2000 words)
\section{Implementation}
\label{sec:chapter-label-implementation}

[Prose explanation before the listing.]

\begin{lstlisting}[
  style=python,
  caption={[Description of what this code does.]},
  label={lst:chapter-label-basic},
]
def example():
    """Docstring."""
    return "result"
\end{lstlisting}

[Prose after listing: \autoref{lst:chapter-label-basic} demonstrates...]

%% SECTION 7: EXAMPLES (2–5 worked examples)
\section{Examples}
\label{sec:chapter-label-examples}

\subsubsection{Example: [Scenario title]}

\textbf{Problem}: [One-paragraph problem statement.]

\begin{lstlisting}[
  style=python,
  caption={Example: [scenario title].},
  label={lst:chapter-label-ex1},
]
# Solution code
\end{lstlisting}

\textbf{Result}: [What the reader should observe.]

%% SECTION 8: COMMON MISTAKES (3–7 items)
\section{Common mistakes}
\label{sec:chapter-label-mistakes}

\subsubsection{Mistake 1: [Short name]}
\textbf{What happens}: [The incorrect usage.]
\textbf{Why it is wrong}: [Root cause.]
\textbf{The correct approach}: [What to do instead.]

%% SECTION 9: BEST PRACTICES (3–7 items)
\section{Best practices}
\label{sec:chapter-label-practices}

\subsubsection{Practice 1: [Short name]}
[What to do, why it matters, how to apply it.]

%% SECTION 13: EXERCISES (3–10)
\section{Exercises}
\label{sec:chapter-label-exercises}

\textbf{Exercise 1} (\(\bullet\) Conceptual)

[Conceptual question testing definitions and principles.]

\textbf{Exercise 2} (\(\bullet\bullet\) Applied)

[Task requiring writing or modifying code.]

\textbf{Exercise 3} (\(\bullet\bullet\bullet\) Advanced)

[Open-ended or research-level challenge.]

%% SECTION 14: SUMMARY (150–300 words)
\section{Summary}
\label{sec:chapter-label-summary}

[Synthesize the chapter. Name the 3–5 most important concepts. No new information.]

%% SECTION 15: FURTHER READING
\section{Further reading}
\label{sec:chapter-label-further}

\textbf{Foundational papers}
\begin{itemize}
  \item \fullcite{Author2024} — [One sentence: why to read this.]
\end{itemize}

\textbf{Official documentation}
\begin{itemize}
  \item [Official source name.] \url{https://example.com} — [One sentence annotation.]
\end{itemize}

\textbf{Books}
\begin{itemize}
  \item \fullcite{Kleppmann2017} — [One sentence annotation.]
\end{itemize}
